% Isolated characteristic-two geometry. No actual surface realization is asserted.
\subsection{The exceptional prime-five polar is a supersingular secant quintic}

Work over an algebraically closed field of characteristic two. In the
prime-five tensor take \(c_1=c_2=1\), and put
\[
 \sigma_i(Y)=Y_i^2+Y_{i-1}Y_{i+1}+Y_{i-2}Y_{i+2},\qquad
 Q(Y)=\sum_{i\in\mathbf Z/5}Y_i\sigma_i(Y)^2.
 \tag{SQ5.1}
\]
All indices in this subsection are modulo five.

\begin{lemma}[A direct characteristic-two model of the Pfaffian curve]
\label{lem:prime-five-supersingular-Pfaffian-curve}
\status{Candidate; complete independent proof/code audits pending.}
The projective scheme cut out by the five quadrics \(\sigma_i\)
is the elliptic normal quintic obtained from
\[
                     E:\quad y^2+y=x^3+x
 \tag{SQ5.2}
\]
by the complete system \(|5O|\), using the following ordered basis:
\[
 \begin{split}
 w_0&=1+x+x^2,&w_1&=1+x+y+x^2+xy,\\
 w_2&=1+xy,&w_3&=1+x+xy,\\
 w_4&=y+x^2+xy.
 \end{split}                                      \tag{SQ5.3}
\]
In particular this curve is smooth and supersingular.
\end{lemma}
\begin{proof}
The smooth projective Weierstrass curve (SQ5.2) has pole orders
\(2,3\) for \(x,y\) at \(O\). Thus
\(1,x,y,x^2,xy\) is a basis of \(H^0(E,5O)\).
Every degree-five line on an elliptic curve separates all
length-two subschemes: the obstruction is \(H^1\) of a line
of degree three, which vanishes. Hence its complete system is
an embedding. The linear change (SQ5.3) has determinant one
over \(\mathbf F_2\).

For an exact scheme check, use coordinates
\((Z_0,Z_1,Z_2,Z_3,Z_4)=(1,x,y,x^2,xy)\). Its five quadrics are
\[
 \begin{split}
 q_0&=Z_1^2+Z_0Z_3,\\
 q_1&=Z_1Z_2+Z_0Z_4,\\
 q_2&=Z_2^2+Z_0Z_2+Z_1Z_3+Z_0Z_1,\\
 q_3&=Z_1Z_4+Z_2Z_3,\\
 q_4&=Z_2Z_4+Z_0Z_4+Z_3^2+Z_0Z_3.
 \end{split}                                      \tag{SQ5.4}
\]
On \(Z_0=1\), the first two eliminate \(Z_3,Z_4\), and
the remaining equations give exactly (SQ5.2). At \(Z_0=0\),
the support of their zero scheme is the single point
\([0:0:0:0:1]\). On \(Z_4=1\), write \(u=Z_3,v=Z_2\).
The equations eliminate
\(Z_1=uv\), \(Z_0=uv^2\), and the remaining ideal is generated by
\[
                    v+u^2+uv^2+u^2v^2.            \tag{SQ5.5}
\]
Its derivative in \(v\) is one. These two affine checks cover
the projective zero scheme and show that it is exactly the smooth
completion of (SQ5.2), with no additional or embedded projective
component.

Substituting (SQ5.3) into the five \(\sigma_i\) gives, in order,
\[
 \begin{split}
 q_0+q_1+q_2,quad&q_0+q_1+q_2+q_4,quad q_0+q_3+q_4,\\
 &q_1+q_3+q_4,quad q_0+q_2+q_4.
 \end{split}                                      \tag{SQ5.6}
\]
This change of the five equations also has determinant one over
\(\mathbf F_2\). It proves the asserted equality of projective
schemes. Finally the invariant regular differential of (SQ5.2)
is \(dx\); its Cartier image is zero, since it is exact.
Thus the elliptic curve has two-rank zero and is supersingular.
No characteristic-zero Hesse discriminant is used.
\end{proof}

\begin{proposition}[The polar quintic is exactly the secant hypersurface]
\label{prop:prime-five-supersingular-secant-polar}
\status{Candidate; complete independent proof/code audits pending.}
The polynomial \(Q\) is irreducible and reduced, and its hypersurface
is precisely \(\operatorname{Sec}(E)\subset\mathbf P^4\) for the
embedded curve above. Its singular support is \(E\); its singular
scheme is defined by \((\sigma_0^2,\ldots,\sigma_4^2)\), and is
not asserted to be the reduced curve.
The same conclusions hold, up to projective coordinates, at every
exceptional ratio \(c_1=1,c_2=t\) with \(t^5=1\).
\end{proposition}
\begin{proof}
For two points \(P,R\) of the curve write \(p,w\) for their
coordinate vectors and let
\(b_i(p,w)=\sigma_i(p+w)+\sigma_i(p)+\sigma_i(w)\).
The following two identities hold in the coordinate ring of
\(E\times E\):
\[
                   \sum_i p_i b_i(p,w)^2=0,
          \qquad   \sum_i w_i b_i(p,w)^2=0.         \tag{SQ5.7}
\]
They can be verified directly by substituting (SQ5.3) in the
two independent copies of (SQ5.2). The retained exact program
does so in
\(\mathbf F_2[x,y,u,v]/(y^2+y+x^3+x,\,v^2+v+u^3+u)\),
reducing every monomial; it records both zero remainders.
Thus these are universal identities on the two-point scheme,
not tests at finitely many points.

Since \(\sigma(p)=\sigma(w)=0\), along their line one has
\(\sigma(\lambda p+\mu w)=\lambda\mu b(p,w)\).
The two coefficients of
\(Q(\lambda p+\mu w)\), at \(\lambda^3\mu^2\) and
\(\lambda^2\mu^3\), vanish by (SQ5.7). Hence \(Q\) vanishes
on every secant and therefore on their closure.

Here is a characteristic-independent intersection calculation of
the degree of that closure. On \(S=\operatorname{Sym}^2E\),
let \(T\) be the rank-two tautological bundle whose fibre at a
length-two divisor \(D\) is \(H^0(E,\mathcal O_E(5O)|_D)\).
Evaluation onto it is surjective by elliptic Riemann--Roch.
The resulting \(\mathbf P(T^*)\to\mathbf P^4\) parametrizes
the secant lines. Four distinct points of \(E\) impose four
independent conditions on \(H^0(5O)\), because the remaining
line has degree one. Consequently two secants with disjoint
endpoints cannot meet; secants with a shared endpoint meet only
there. Away from \(E\) the general point therefore has a
unique unordered pair of endpoints. The tangent lines at two
distinct points span a projective three-space, because the
length-four divisor \(2P+2R\) likewise imposes four independent
conditions. The differential of the secant parametrization has
rank three at a general point. This proves that the image is an
irreducible hypersurface and that the parametrization is
birational, without using generic injectivity as a substitute
for separability.

Pull \(T\) back under the degree-two finite map
\(\rho:E\times E\to S\). Evaluation at the first point gives
an exact sequence of line bundles
\[
 0\longrightarrow p_2^*\mathcal O_E(5O)(-\Delta)
 \longrightarrow\rho^*T
 \longrightarrow p_1^*\mathcal O_E(5O)\longrightarrow0.
 \tag{SQ5.8}
\]
It remains valid on the diagonal, where it is the usual
value-and-first-jet sequence. If \(f_1,f_2\) are the two
fibre classes, then \(f_i^2=\Delta^2=0\),
\(f_1f_2=f_i\Delta=1\). Thus
\[
 \int_{E\times E}c_1(\rho^*T)^2=30,
 \qquad \int_{E\times E}c_2(\rho^*T)=20.
 \tag{SQ5.9}
\]
The projective-bundle formula and the degree of \(\rho\) give
\(\int_{\mathbf P(T^*)}c_1(\mathcal O(1))^3=(30-20)/2=5\).
These are integer intersection degrees; no division by two in
the ground field is involved. Birationality then gives
\(\deg\operatorname{Sec}(E)=5\).
Since \(Q\) has degree five and is nonzero (its \(Y_0^5\)
coefficient is one), it is the reduced irreducible defining
equation. Its ordinary derivatives are \(\sigma_i^2\).
Euler's identity with odd degree five gives
\(Q=\sum Y_i\sigma_i^2\), so the singular ideal is exactly
the one stated, with reduced support determined by the lemma.

Finally rescale a Schr\"odinger basis by
\(e_i\mapsto\zeta^{a i^2}e_i\), where \(\zeta^5=1\).
The tensor coefficient at distance \(h\) is multiplied by
\(\zeta^{a h^2/2}\). Hence the ratio \(c_2/c_1\) is
multiplied by \(\zeta^{4a}\). These phases carry each of the
five roots \(t^5=1\) to one, up to a nonzero common scalar.
The corresponding source and dual target coordinate changes
preserve the incidence expression \(y(s^{[2]})\), so carry
the Pfaffian and polar geometry to the model already proved.
\end{proof}

\begin{remark}[Certificate and geometric scope]
\label{rem:prime-five-supersingular-secant-scope}
\status{Candidate scope; no actual Cartier realization asserted.}
The reconstructive program is
\path{code/monodromy/post_CP36_prime5_ss_quintic_certificate.py}.
It passed at 06:06:37.270415 UTC on 20 November 2026 in
\(0.002333582000574097\) seconds, checking both changes of
basis, all five transformed quadrics, the complete smooth
infinity chart and both universal secant identities. The earlier
binary cyclic-basis search and the unsuccessful comparison with
five Gaussian products of linear forms are retained separately.
The latter comparison does not assert irreducibility; that
assertion follows from the geometric degree calculation above.

This identifies the abstract exceptional polar in characteristic
two. It neither asserts that a particular abelian-surface Cartier
tensor has these coefficients nor excludes that possibility.
In particular a ruled or contracted image is not automatically
incompatible with a globally generated type-\((1,5)\) polarization.
No monodromy group or common wall is deduced here.
\end{remark}
